\documentclass[a4paper]{article}

% Set margins
\usepackage[hmargin=2.5cm, vmargin=3cm]{geometry}

\frenchspacing

% Language packages
\usepackage[utf8]{inputenc}
\usepackage[T1]{fontenc}
%\usepackage[magyar]{babel}

% AMS
\usepackage{amssymb,amsmath}

% Graphic packages
\usepackage{graphicx}

% Colors
\usepackage{color}
\usepackage[usenames,dvipsnames]{xcolor}

% Listings
\usepackage{listings}

% Question
\newenvironment{question}[1]
{\noindent\textcolor{OliveGreen}{$\circ$ \textit{#1}}

\smallskip

\color{Gray}

}{\bigskip}

% Task
\newenvironment{task}[1]
{\noindent\textcolor{RoyalBlue}{$\circ$ \textit{#1}}

\smallskip

\color{Gray}

}{\bigskip}

% Notification
\newenvironment{notification}[1]
{\noindent\textcolor{Peach}{$\circ$ \textit{#1}}

\smallskip

\color{Gray}

}{\bigskip}

% Problem
\newenvironment{problem}[1]
{\noindent\textcolor{OrangeRed}{$\circ$ \textit{#1}}

\smallskip

\color{Gray}

}{\bigskip}

%===============
\begin{document}
%===============

\definecolor{codebg}{rgb}{0.9,0.9,0.9}
\lstset{backgroundcolor=\color{codebg}}

\begin{center}
    \Huge \textbf{Programtervezési ismeretek}
    
	\bigskip
	
	\Large Feladatok gyakorláshoz
\end{center}

\bigskip

%=====================
\section*{1. Halmazok, műveletek}
%---------------------

\begin{enumerate}

\item

Tekintsük az $A = \{\alpha, \beta, \gamma, \zeta\}$ és $B = \{igaz, hamis\}$ halmazokat!
\begin{itemize}
	\item Írjuk fel az $A \times A, A \times B, B \times A, B \times B$ Déscartes szorzatokat!
	\item Írjuk fel a $2^A$, $2^B$ halmazokat!
	\item Írjuk fel az $A^2 \times B$ halmazt!
	\item Mennyi eleme lesz az $A^2 \cup (B \times A)$ halmaznak?
	\item Teljesül-e a következő egyenlőtlenség?
	$$
	|(A \times B) \cup (B^2 \times A^2)| < |A^3 \cup B^3|
	$$
\end{itemize}

\item

Számítsuk ki a következő halmazok elempárjainak szorzatát $a \cdot b$ formában orosz-paraszt módszerrel!
$$
S = \{(48, 13), (59, 108), (1024, 127), (50, 128), (123, 456)\}
$$
\begin{itemize}
	\item Ellenőrízzük kézzel kiszámítva a szorzatot!
	\item Számítsuk ki $b \cdot a$ formában is!
\end{itemize}

\item

Számítsuk ki az
$$
\lfloor a \rfloor, \lceil a \rceil,
\lfloor b \rfloor, \lceil b \rceil,
\left\lfloor \dfrac{a}{b} \right\rfloor, \left\lfloor \dfrac{b}{a} \right\rfloor, 
\left\lceil \dfrac{a}{b} \right\rceil, \left\lceil \dfrac{b}{a} \right\rceil, 
$$
$$
\text{Round}\left( \dfrac{a}{b} \right),
\text{Round}\left( \dfrac{b}{a} \right),
\left\{ \dfrac{a}{b} \right\},
\left\{ \dfrac{b}{a} \right\}
$$
kifejezések értékét, az összes lehetséges $a$-ra és $b$-re!
$$
A = \{3, -2, 8, 1, 0.2, \pi\}, \quad
B = \{5, 7, -4, 0.9, -e\}
$$

\item

Írjuk fel a következő halmazok elemeit!
$$
P = \{a \text{ mod } b | a \in \{5, 0, -2, 9\}, b \in \{-1, 4, 21, 0, 98\}\}
$$
$$
Q = \{a \text{ div } b | a \in \{50, 41, -28, 11\}, b \in \{-4, 72, 18\}\}
$$

\end{enumerate}

%=====================
\section*{Számrendszeri átváltás}
%---------------------

\begin{enumerate}

\item

Egész szám adatstruktúra

$$
c_n, c_{n-1}, \ldots, c_1, c_0 \quad 0 \leq c_k < b, b \geq 2
$$

$$
x = c_n \cdot b^n + c_{n-1} \cdot b^{n-1} + \cdots + c_1 \cdot b^1 + c_0 \cdot b^0 = \sum_{i=1}^n c_n \cdot b^n
$$

\begin{itemize}
	\item Írjuk fel a 5291-et 2, 3, 7, 8, 10, 11, 16-os számrendszerben!
	\item Ellenőrízzük helyiértékesen és Horner sémával!
\end{itemize}

\item

Mennyi számjegyből fog állni a 12345678 szám 2, 3, 6, 10, 15, 16-os számrendszerben?

$$
n + 1 = \left\lfloor \log_b x \right\rfloor + 1
$$

\item

Negatív egészek ábrázolása, kettes komplemens

Mennyi byte szükséges a -11, -923, -30000, -300000 számok ábrázolásához?

\item

Törtek átírása
\begin{itemize}
	\item Írjuk át a 0.527, -0.32 és -0.64 értékeket kettes és tizenhatos számrendszerbe!
	\item Írjuk át a $0.110110110_2$, $0.21022_3$ és $0.AFFE_{16}$ értékeket tizes számrendszerbe!
\end{itemize}

Váltsuk át a 32412 és -1111 számokat kettes számrendszerbe!
\begin{itemize}
	\item Mennyi bitből fog állni?
	\item Adjuk össze őket kettes számrendszerben!
	\item Az összeget számoljuk vissza 10-es számrendszerbe, és ellenőrízzük!
\end{itemize}

\item

Írjuk fel kettes számrendszerbe a következő értékeket!
$$
0.42578125_{10}, \quad
0.70312_{10}, \quad
0.74_{10}, \quad
0.14062_{10}, \quad
0.15_{10}, \quad
0.ABCD_{16}
$$
Ellenőrízzük helyiértékesen és Horner sémával is!

\item

Ha egy szám 83 jegyű kettes számrendszerben, akkor hány jegyű lesz 16, 12 és 10-es számrendszerben?

\medskip

Ha egy szám 20 jegyű 16-os számrendszerben, akkor hágy jegyű lesz 10-esben, 8-asban és kettesben?

\item

Mennyi byte-on lehet ábrázolni a TAJ számot?

\end{enumerate}

%=====================
\section*{Lebegőpontos számábrázolás}
%---------------------

\begin{enumerate}

\item

Ábrázoljuk egyszeres lebegőpontos számábrázolással a következő értékeket!
$$
4152.6374, \quad
-0.49, \quad
5592.134, \quad
0.83, \quad
99.77
$$

\item

Számítsuk ki, hogy milyen értékeket ábrázolnak a következő byte-ok egyszeres lebegőpontos ábrázolást feltételezve!
$$
9FC40000, \quad FF88C000, \quad 14130000, \quad C5C50000, \quad 87654321
$$

\end{enumerate}

%=====================
\section*{Logikai műveletek, kapuáramkörök}
%---------------------

\begin{enumerate}

\item

Fejezzük ki a bináris logikai műveleteket az $\wedge$, $\vee$ és negáció műveletekkel!

\item

Írjuk fel a 5 bemenetes többségi szavazás diszjunktív normál formáját, és rajzoljuk fel a kapuáramkört!

\item

Adjuk meg a végtelen értékeket egyszeres pontossággal!

\item

Milyen érték a 7F801110?

\item

Tervezzük meg egy 7 bemenetes automatának a logikai kapuáramkörét, amely a maximumot adja vissza!

\item

Tervezzünk kapuáramkört $\wedge$, $\vee$ és negáció kapukból az $u = f(x, y, z) = \overline{(x \oplus y)} \rightarrow \overline{z}$ logikai függvényhez!

\item

Vizsgáljuk meg, hogy teljesülnek-e a következő azonosságok!

$x \wedge (y \rightarrow z) \equiv (x \rightarrow z) \wedge (x \rightarrow z)$

$x \oplus (y \vee z) \equiv (x \oplus y) \vee (x \oplus z)$

$(c \oplus d) \rightarrow (e \leftrightarrow d) \equiv c \vee (e \wedge d) \vee (\overline{e} \wedge \overline{d})$

$a = (a \rightarrow b)|(b \rightarrow \overline{a})$

$x \wedge (x \vee \overline{y}) \wedge (x \vee \overline{z}) \equiv ((\overline{x} 
\downarrow y) \wedge \overline{z}) \vee x$

\item

Lássuk be, hogy a Sheffer vonás nem asszociatív!

\item

Tervezzünk kapuáramkört az $\wedge$, $\vee$ és negáció műveletekkel az $u = f(x, y, z) = (x \oplus (y \rightarrow z)) \wedge (x \rightarrow y)$ logikai függvényhez!

\item

Tervezzünk 8 bemenetes paritás ellenörző automatát!

\item

Igazoljuk a következő egyenlőségeket!

$\overline{\overline{X \cap Y} \cup X} = \emptyset$

$\overline{A \cap \overline{B \cup A} \cap B} = \overline{\emptyset}$

$
(E \cap F \cap \overline{G}) \cup (E \cap \overline{F} \cap G)
=
E \setminus (G \triangle \overline{F})
$

\item

Mennyi olyan halmazpár van, amelyeknek a metszete $\{\mu, \sigma\}$ uniója pedig $\{\xi, \eta\}$?

\item

Mik lesznek a következő halmazok elemei?

\begin{itemize}
    \item $\{x \in \mathbb{R} : x^3 - 2x^2 + 8 = 0\}$
    \item $\{x \in \mathbb{R} : 3\log x = \log x^3\}$
    \item $\{x \in \mathbb{R} : 5 < |x| \leq 10\}$
    \item $\{(x, y) \in \mathbb{R}^2 : x^2 + y^2 = 16 \text{ és } y \geq 0\}$
    \item $\{(x, y) \in \mathbb{R}^2 : |x| < y \leq 4\}$
\end{itemize}

\item

Tekintsük az $A = (-15; 23]$ és a $B = [15; 21)$ intervallumokat. Határozzuk meg a következő kifejezések értékét:
$A \cup B, A \cap B, A \setminus B, B \setminus A, A \bigtriangleup B, A \times B, 2^A, 2^B$.

\item

Milyen elemek tartoznak a következő halmazokhoz?

\begin{itemize}
    \item $\{(a, b) \in \mathbb{N}^2 | a \mod b = 8, a + b = 50\}$
    \item $\{(a, b) \in \mathbb{N}^2 | a \text{ div } b = 4, a \cdot b < 100\}$
    \item $\{(x, y) \in \mathbb{R}^2 : \text{Round}(x) + \text{Round}(y) = 5\}$
    \item $\{(x, y) \in \mathbb{R}^2 : \{3 \cdot x\} > 2 \cdot \{y\}\}$
\end{itemize}

\item

Írjuk fel a $2^{2^{\{a, b, c\}} \setminus \{(x, y)|x \in\{a, b\},  y \in\{b, c\}\}}$ halmaz elemeit!

\item

Ha $|H| = 122, A \subset H, B \subset H, |A| = 79, |B| = 61$ akkor mennyi lesz $|A \cup B|$ és $|A \cap B|$?

\end{enumerate}

%=====================
\section*{Karak adattípus, Unicode}
%---------------------

\begin{enumerate}

\item

Alakítsuk át a következő unicode szimbólumokat UTF-8 kódolásúra!

U+0052, \quad U+09AE, \quad U+10EAC8, \quad U+88FF

\end{enumerate}

%=====================
\section*{Sor- és oszlopfolytonos tárolási mód}
%---------------------

\begin{enumerate}

\medskip

\item

Adott egy tömb a memóriában, amelynek elemei 24 bájtosak. A 92-edik elem a tömbelem 10. bájtjának 3. bitje hanyadik bitje lesz a tömbnek?
\begin{itemize}
	\item Vizsgáljuk külön a 0 és 1 kezdőindexek eseteit!
	\item Hogyan tudnánk felírni általánosan, ha $m$ bájtos a tömb, $i$-edik tömbelem $j$-edik bájtjának $k$-adik bitjéről van szó?
\end{itemize}

\end{enumerate}

%=====================
\section*{Procedúrák, rekurzív algoritmusok}
%---------------------

\begin{enumerate}

\item

Konkatenáljuk össze egy karaktereket tartalmazó vektor első 10 elemét!
\begin{itemize}
	\item Írjuk fel a procedúra be- és kimeneteit!
	\item Írjuk fel a procedúra és a főprogram működését lépésenként!
	\item Vizsgáljuk a futás közben a verem állapotát!
\end{itemize}

\item

Készítsünk rekurzív procedúrát az előző feladatra!
\begin{itemize}
	\item Hogyan változik a verem állapota futás közben?
	\item Rajzoljuk fel a rekurzív hívási fát!
	\item Mennyi rekurzív hívás történt, és mennyi volt a maximális veremmélység?
	\item Milyen előnyei/hátrányai vannak a rekurzív megvalósításnak?
\end{itemize}

\item

Számítsuk ki az $\binom{10}{9}$ értékét!

$$
\binom{n}{k} = \binom{n-1}{k-1} + \binom{n-1}{k}
$$

\item Számítsuk ki a $\sqrt{111}$ értékét az iteratív képletnek megfelelően $\varepsilon = 0.0001$ pontosságig!
$$
x_{k+1} = \dfrac{1}{2} \cdot \left( x_k + \dfrac{s}{x_k} \right)
$$

\item

Számítsuk ki a Div\_Sum(2000, @$x$) hívás esetén az $x$ értékét!

Div\_Sum(a, @r) \\
// Input: $a \in \mathbb{Z}$ \\
// Output: $r \in \mathbb{Z}$ \\
IF a < 20 \\
\indent \quad THEN $r \leftarrow 0$ \\
\indent \quad ELSE CALL Div\_Sum($a$ DIV 5, @$r$) \\
\indent \quad \quad \quad $r \leftarrow r$ + ($a$ DIV 3) + ($a$ DIV 4) \\
RETURN(r)

\item

Írjunk egy procedúrát, amely egy tömbhöz egy olyan tömböt számol ki, melynek $i$-edik elemén az eredeti tömb $i$-edik előtti elemeinek összege szerepel!

\item Írjunk fel egy rekurzív függvényt a faktoriális kiszámítására!
\begin{itemize}
	\item Számítsuk ki az $8!$ értékét!
	\item Mennyi rekurzív hívás volt a számítás során?
	\item Mekkora volt a maximális mélysége a hívási veremnek?
	\item Rajzoljuk fel a rekurzív hívási fát!
\end{itemize}

\item A bináris keresés algoritmusával keressük meg a $200$ értéket a következő tömbben!

$$
A = [50, 100, 120, 140, 140, 200, 280, 432]
$$

\item

Készítsünk egy procedúrát, amely megállapítja, hogy egy vektor mely elemei esnek egy adott intervallumba!

\item

Ellenőrízzük, hogy egy tömb elemei nemnövekvő sorrendben vannak-e!

\item

Készítsünk \textit{egy} procedúrát, amely kiszámítja két vektor különbségét és diadikus szorzatát!

\item

Írassuk ki egy vektor értékeinek mediántól való eltérését!

\item

Ellenőrízzük, hogy két vektorban elemei között nincsenek azonos értékűek!

\item

Hogyan tudnánk bekérni, tárolni és kiíratni egy felsőháromszögmátrix értékeit?

\item

Vizsgáljuk meg, hogy egy mátrix diagonális-e!

\item

Egy egész értékeket tároló vektorban keressünk olyan számhármasokat, amelyekre teljesül, hogy $2x - 3y = z^3$.

\item

Síkbeli pontok adatait tároljuk egy tömbben $(x, y)$ rendezett párok formájában folytonosan.

$$ [
\begin{array}{ccccccccc}
x_0 & y_0 & x_1 & y_1 & \ldots & x_{n-1} & y_{n-1} & x_n & y_n
\end{array}
] $$

\item

Írassuk ki a koordináták mellé az origótól való távolságukat!

\item

Forgassuk el az összes pontot az origó körül egy adott szöggel!

\item

Írassuk ki azon pontok koordinátáit, amelyek nem esnek az origó $r$ sugarú környezetébe!

\item

Keressük meg azon pontpárokat, amelyek távolsága egy adott $\delta$ értéknél kisebb!

\item

Keressük meg azon pontnégyeseket, amelyek egy körre esnek!

\item

Általánosan hogy kereshetnénk legkisebb területű szabályos sokszögeket, amelyeknek csúcsai az adott pontok?

\item

Készítsünk egy procedúrát, amely egy szöveges értékről megállapítja, hogy az lehet-e telefonszám!

\item

Távolítsuk el egy szövegből az ékezetes karaktereket!

\end{enumerate}

%======================
\section*{Struktúrált programozás}
%----------------------

\begin{enumerate}

\item

Rajzoljuk fel a következő élhalmazokhoz a programgráfokat!

teljes
\begin{gather*}
P = \{
(START, A),
(A, B/i), (A, F/h),
(B, E), (E, G), \\
(F, C), (C, D/i), (C, G/h),
(G, H), (H, STOP)
\}
\end{gather*}

tömör
\begin{gather*}
P = \{
START \rightarrow P,
P \rightarrow (Q, R),
Q \rightarrow (T, S),
R \rightarrow (S, U),
S \rightarrow (V, W), \\
T \rightarrow X,
V \rightarrow X,
U \rightarrow Y,
W \rightarrow Y,
X \rightarrow STOP,
Y \rightarrow STOP\}
\end{gather*}

Vizsgáljuk meg, hogy a program, és ha nem, akkor rajzoljuk fel az ekvivalens struktúrált program gráfját!

\begin{itemize}
	\item írjuk fel a struktúrált program formuláját!
	\item Írjuk fel az eredeti és a struktúrált program pszeudó kódját!
	\item Számítsuk ki a ciklikus bonyolultságot!
	\item Rajzoljuk fel a struktogramot!
\end{itemize}

\item

Rajzoljuk fel a programgráfot a pszeudókód alapján!

\begin{verbatim}
1 A
2    IF B
3        THEN GOTO 6
4        ELSE GOTO 1
5 D
6 IF C
7    THEN GOTO 2
\end{verbatim}

\end{enumerate}

\end{document}

\item

A pszeudókódok alapján rajzoljuk fel a programgráfot, majd készítsük el az ekvivalens struktúrált programokat.

\begin{verbatim}
1 IF A
2   THEN B
3       IF C
4           THEN F
5           ELSE GOTO 8
6   ELSE
7       WHILE D DO
8           E
9 G
\end{verbatim}

\begin{verbatim}
1 IF A
2   THEN
3       WHILE B DO
4           C
5       GOTO 10
6   ELSE
7       IF D
8           THEN
9               DO E
10                  WHILE F
11          ELSE G
\end{verbatim}

\begin{verbatim}
1 IF A
2   THEN
3       B
4       IF D
5           THEN GOTO 7
6   ELSE
7       C
8       IF E
9           THEN GOTO 3
\end{verbatim}

\item

Keressük meg egy valós értékeket tartalmazó tömbben egy procedúrával a paraméterként kapott $z$ értékhez a legközelebbi elemet!

\item

Válogassuk szét egy valós tömb elemeit negatív és nemnegatív értékekre!

\item

Töröltjük egy tömbből a 0 értékű elemeket! \textit{(a. segédtömbbel, b. segédtömb nélkül)}

\item

Írjunk egy procedúrát, amely 2, azonos méretű valós értékeket tároló tömbhöz visszaadja az adott indexeken lévő minimális értékeket.

pl.: $x = [5, 8, -1], y = [2, 1, 3] \Rightarrow z = [2, 1, -2]$

\item

Számítsuk ki két vektor különbségének négyzetösszegét!

\end{enumerate}

%======================
\section*{13. gyakorlat}
%----------------------

\textit{COCOMO modell}

\begin{itemize}
    \item PM: Person Month
    \item $T_d$: Development Time
    \item KDSI: Kilo Delivered Source Index
\end{itemize}

$$
PM = 2.4 \cdot KDSI^{1.05},
\quad
T_d = 2.5 \cdot \sqrt[3]{PM},
\quad
\text{létszám} = \dfrac{PM}{T_d},
\quad
\text{Brook limit}: T_d \cdot 0.75
$$

\begin{enumerate}

\item Egy 100000 soros alkalmazást szeretnénk kifejleszteni. Mennyi emberre és időre van szükség?

\item 20 programozóval fél év alatt mekkora alkalmazást lehet kifejleszteni?

\item Egy 300000 soros program esetében lehetőségünk van egy függvénykönyvtár megvásárlására, amely 120000 sor kiváltására képes.

\medskip

Maximum mennyit érdemes fizetnünk ezért, ha a programozók havi fizetése 240000 Ft?

\item Egy 50000 soros alkalmazást szeretnénk minél gyorsabban elkészíteni. Mennyi plusz embert lehet még bevonni, hogy a határidő tartása ne legyen rizikós?

\item Egy 1000000 soros alkalmazást 25 emberrel szeretnénk kifejleszteni. Az alkalmazást várhatóan 5 millió forintént fog elkelni.

\medskip

Legfeljebb mennyi lehet így a programozók havi fizetése?

\item

Írjuk át iteratív alakra a következő rekurzív algoritmust!

Parity(x, @y) \\
// Input: $x \in \mathbb{Z}$ \\
// Output: $y \in \mathbb{Z}$ \\
IF x = 0 \\
\indent \quad THEN y $\leftarrow$ 1 \\
\indent \quad ELSE IF x mod 2 = 0 \\
\indent \quad \quad THEN CALL Parity(x-1, y) \\
\indent \quad \quad \quad \quad \quad y $\leftarrow$ y + x \\
\indent \quad \quad ELSE CALL Parity(x-1, y) \\
\indent \quad \quad \quad \quad \quad y $\leftarrow$ y $\cdot$ x \\
RETURN(y)

\item

Írjuk át iteratív alakra a következő rekurzív algoritmust!

ADDER(@X, i, @a, @b) \\
// Input: $X \in \mathbb{R}, i \in \mathbb{N}$ \\
// Output: $a, b \in \mathbb{R}$ \\
IF i > 0 \\
\indent \quad THEN \\
\indent \quad \quad ADDER(X, i - 1, a, b)
\indent \quad \quad IF $x_i$ > 0 \\
\indent \quad \quad \quad THEN a $\leftarrow$ a + xi \\
\indent \quad \quad \quad ELSE b $\leftarrow$ b + xi \\
\indent \quad ELSE \\
\indent \quad \quad a $\leftarrow$ 0 \\
\indent \quad \quad b $\leftarrow$ 0 \\
RETURN(a, b)

\item

Írjuk át iteratív alakra a következő rekurzív algoritmust!

TRNG(x, @y) \\
// Input: $x \in \mathbb{R}$ \\
// Output: $y \in \mathbb{R}$ \\
IF x > 0 \\
\indent \quad THEN FOR i $\leftarrow$ 1 TO X DO \\
\indent \quad \quad \quad \quad \quad CALL TRNG(x-1, y) \\
\indent \quad \quad \quad \quad \quad INC(y) \\
\indent \quad ELSE y $\leftarrow$ 0 \\
RETURN(y)

\item

Készítsünk egy algoritmust, amely úgy osztályozza a tömbben tárolt egészeket, hogy a párosakat a tömb elejére, a páratlanokat a tömb végére helyezi.

\item

Készítsünk algoritmust, amely $k$ szerinti maradékosztályok szerint rendezi egy tömb elemeit, vagyis az egyes szakaszok értékeit $k$-val osztva $0, 1, \ldots k-1$ maradékot adnak.

\medskip

Például 4 esetén:

$$
[1, 5, 9, 7, 4, 3, 1, 2, 8, 6, 5] \quad \rightarrow \quad [4, 8, 1, 5, 9, 1, 5, 2, 6, 7, 3]
$$

\item

Írassuk ki egy vektor elemeinek összes lehetséges permutációját!
\begin{itemize}
    \item[a.] Rekurzív algoritmussal
    \item[b.] Iteratív algoritmussal
\end{itemize}

\item

Számítsuk ki, hogy egy vektorban melyik értékből mennyi van.
\begin{itemize}
    \item[a.] Természetes számokat tároló vektor előre rögzített maximális értékkel
    \item[b.] Tetszőleges elemtípusra vektorok vektorával
\end{itemize}

\end{enumerate}

\end{document}
